\chapter{Discussion and Conclusions}
\label{chap:discussion}

\section{Answers to the research questions}

\subsection{RQ1: experimental reliability}

The observed RT-DETR variability is large enough to change single-run model
rankings. Nominally identical native-CUDA executions diverge after the first
backward pass even when initialization, input, and initial loss agree. The
transferred COCO person head reduces amplification of those perturbations but
does not remove them, and DCNv2 prevents strict deterministic FAM training in
the studied environment.

Reliable comparison consequently requires more than recording a seed. The
historical locked protocol combines paired seed values, isolated processes, a
fixed training horizon, epoch-10 checkpoints, identical pretrained-head
transfer, training-only Modal Dropout, and distributional reporting. It turns
the single-run SSJ result into a repeatable selection of standard FAM over the
Base RT-DETR configuration.

The validation-led campaign additionally shows that the checkpoint-selection rule
materially affects the measured result. Holding out the complete FHL
0401/0402 video and selecting \texttt{best} within a fixed ten-epoch budget
improves MtErie over \texttt{latest} in all five FAM Stage-A seeds. This is
checkpoint selection, not early stopping: every run still completes ten
epochs. Separating validation-led Stage A from full-data Stage B also prevents
an apparent gain from being created by comparing different optimization
datasets or code revisions.

The methodological answer to RQ1 is therefore that experimental design is part
of the result. Paired replication, a complete-video development split,
predeclared promotion rules, and a configuration-matched full-data confirmation
are all needed before replacing the baseline. The selected split is still
internal, and MtErie remains previously consulted rather than blind.

The acquisition-level confirmation demonstrates the complementary role
of freezing data inventories, checkpoint hashes, and contrasts before
inference. It provides stronger evidence than another MtErie evaluation, but
does not turn internal WiSARD sequences into an external holdout. The
mixed-consistency, box-guided, RT-DETRv2, and YOLO26 comparisons apply the same
five-seed Stage-A decision structure as the other candidates. No performance
decision is taken on an individual seed: the complete budget is run for every
seed, the validation-selected \texttt{best} is primary, and the final epoch is
a mandatory diagnostic. Promotion requires both a mean paired gain of at least
0.01 and positive gains in at least four of five seeds. The fixed horizon is
ten epochs for RT-DETR and RT-DETRv2 and fifty for YOLO26.

\subsection{RQ2: RT-DETR effectiveness and modality use}

In the historical locked comparison, standard FAM improves Base RT-DETR in
every paired seed and raises mean \mapfifty{} by 0.0700. The
improvement is supported by COCO mAP, mAP@75, mAR, and the recall analysis. IR
feature dropout and SSJ do not provide a reliable increment over FAM; Identity
DCNv2 is highly variable, and Grid Sample is close but not consistently
superior.

The historical native-sensor modality table evaluates VIS, IR, and fusion on
different frame populations and annotations, therefore it cannot support a
direct cross-modality ranking. The post-selection audit holds all of these
variables fixed for the five current-code FAM checkpoints. On the same 708
paired MtErie frames and VIS ground truth, fusion exceeds VIS in 5/5 seeds
by \(0.0373\pm0.0223\), with a 95\% interval of [0.0095, 0.0650]. This provides
the supported evidence that the thermal stream adds information beyond VIS
on the selected internal population.

The very low \(0.0215\pm0.0091\) masked-IR/VIS-ground-truth score is a
different robustness construct. It removes RGB, which supplies the reference
for the asymmetric FAM, but still requests boxes in VIS coordinates from IR
that has only been resized and padded. A post-hoc native-coordinate diagnostic
obtains \(0.5618\pm0.0498\) on the same IR counterparts with native IR
annotations, ruling out a collapsed thermal branch. Replacing native IR dropout
with the paired VIS-coordinate convention fails the frozen seed-40 ablation: it
reduces native-IR performance in that screening run without solving
paired missing-RGB prediction.

The acquisition-level confirmation strengthens the architectural result while
narrowing the modality claim. Historical FAM exceeds Base in all five
seeds on both previously unused acquisitions, with gains of 0.0960 on
Carnation and 0.1554 on FHL. On the matched Stage-B checkpoints, however,
fusion is effectively tied with VIS on Carnation and is worse on FHL. A better
fusion operator does not imply that the thermal input helps uniformly on every
acquisition.

The answer to RQ2 is therefore conditional. FAM consistently improves the
Base RT-DETR reference, and fusion contributes beyond VIS under a
same-frame paired evaluation. Modal Dropout supports native-sensor operation
under its training contract, but it does not make the asymmetric fusion model
robust to every possible missing-reference and coordinate-system intervention.

\subsection{RQ3: mechanism, levels, and improved alignment}

The FAM predicts nonzero transformations at P3 and P4 across all five standard
historical checkpoints. Its offset scale and transformed features vary with
seed, level, and acquisition session. SSJ changes P3 behavior consistently,
increasing offset magnitude and smoothing, even though detection does not
improve. Grid Sample is internally more stable but does not establish that a
pure warp is equivalent to DCNv2.

The diagnostics do not demonstrate physical RGB--IR registration. DCNv2 mixes
geometric sampling, filtering, mask modulation, and channel interaction. The
supported interpretation is a task-oriented geometric-functional
transformation. The P5 seed-41 collapse reinforces this distinction: a
spatially constant internal branch coexists with the best historical FAM
checkpoint because the rest of the detector compensates for it. Smoothly
bounding offsets removes the collapse in five new runs but lowers mean
accuracy, so a mechanistically effective correction is not automatically a
better detector.

The validation-led extension campaign rejects several simple improvement
hypotheses. Adding P2 loses in 5/5 seeds, a matched micro-batch control does not
explain that deficit, global \(800\times800\) input wins only 1/5 comparisons,
and two post-FAM reliability gates fail their promotion rules. These results do
not prove that high-resolution features or reliability modeling are generally
unsuitable; they reject the implemented recipes under the common ten-epoch
protocol.

Zero-offset initialization also fails Stage A: resetting only the initial
sampling offsets reduces best-validation mean \mapfifty{} by 0.0125 and wins
only one of five pairs. The final-epoch mean gain of 0.0118 has three positive
pairs and does not change the decision based on validation-selected best
checkpoints. Both paired confidence intervals include zero. The predictors
learn nonzero, distinct weight rows in every inspected level and checkpoint,
so their initially identical zero rows do not remain identical. This isolates
an initialization intervention from the historical Identity DCNv2 operator,
while providing no evidence of improved physical registration or a reliable
performance advantage under the tested protocol.

RCRA is the only extension to pass Stage A. Its learned coefficient mainly
suppresses the FAM residual at P3 and P4 and reacts to either missing modality
in all five seeds; a three-parameter scalar control cannot reproduce the
validation gain. In matched full-data Stage B, however, RCRA improves mean
\mapfifty{} by 0.0184 but wins only 3/5 seeds, with an interval spanning
[-0.0550, 0.0918]. It therefore fails the frozen stability condition. The
answer to RQ3 is again asymmetric: reliability-conditioned residual selection
is mechanistically active and promising, but its performance gain is not
stable enough to replace standard FAM.

Mixed consistency training and Box-Guided P3 are evaluated against the same
five standard-FAM Stage-A controls as the other RT-DETR candidates. Mixed
consistency reduces best-checkpoint mean \mapfifty{} from
\(0.1646\pm0.0196\) to \(0.1094\pm0.0328\), with negative paired differences
in all five seeds. Box-Guided P3 obtains \(0.1723\pm0.0233\), but its mean
gain of 0.0077 and two positive pairs fail both promotion conditions. Its
final-epoch gain is only 0.0005. Neither recipe is promoted to full-data Stage B.
These results reject the tested detection recipes, not their underlying
mechanistic objectives. Five-seed geometric counterfactuals for box-guided
alignment and isolated-modality evaluations for mixed consistency were not
performed and their outcomes cannot be inferred from the fusion comparison.

The synthetic stress supplies a second limitation. FAM does not exhibit a
universally more favorable signed degradation curve than Base; some FHL
translations improve the unmodified Base reference more. The
supported mechanism remains a task-oriented geometric-functional
transformation, not verified physical registration or invariant tolerance to
known shifts.

\subsection{RQ4: operational relevance}

FAM primarily changes the recall side of the precision--recall trade-off. At
confidence 0.25, it improves mean recall by 0.0990 in all five paired seeds and
reduces the frequency of non-empty frames with at least one missed target. This
is relevant to SAR, where a missed person can have high cost. At a common
confidence threshold, false positives do not decrease consistently, however,
and the 0.01 threshold used for AP collection is not operational because it
produces roughly fifty false positives per image.

The post-hoc recall--FPPI comparison in \Cref{sec:recall-fppi} separates this
threshold-dependent observation from performance at a matched false-positive
budget. At 0.5 FPPI, historical FAM increases mean recall by 0.1102 on MtErie,
0.0981 on Carnation 0025/0026, and 0.1479 on FHL 0407/0408, with positive
differences in all five seeds within each population. This supports a better
recall--false-positive trade-off at the reported budgets, not universal
curve dominance or a validated deployment threshold.

Carnation 0023/0024 exposes a geometric boundary. FAM outperforms Base RT-DETR under the
severe scale mismatch, but adding IR remains worse than using VIS alone in
every FAM checkpoint. The experiment therefore establishes a relative FAM
advantage under mismatch, not correction or mitigation of the underlying
sensor misalignment.

On the two larger previously unused internal acquisitions, FAM's absolute
advantage over Base persists in all ten acquisition--seed comparisons.
The controlled perturbations show, however, that operational robustness cannot
be reduced to a monotone translation or scale curve without calibrated sensor
geometry. These findings motivate acquisition-specific validation rather than
a universal alignment claim.

The performance gain also has a material computational cost: 77.5\% more
parameters, 46.2\% more proxy GFLOPs, 17.9\% more detector latency, and higher
CUDA memory on the measured GPU. These values prevent describing FAM as a free
or generically real-time improvement.

\subsection{RQ5: architectural transferability}

Consistently positive transfer is not established. Deformable DETR provides a
promising exploratory median trend, but its observed ranges overlap and its
protocol is not aligned with the final campaign. Complete DINO base is too weak
under the studied configuration to justify FAM variants. Under the controlled
YOLOv10 protocol, FAM decreases mean VIS+IR performance and wins in only two of
five seeds.

The five-seed Stage-A contrasts distinguish RT-DETRv2 from YOLO26
(\Cref{tab:detector-stage-a-five-seed,tab:detector-stage-a-deltas}). RT-DETRv2
FAM increases best-checkpoint mean \mapfifty{} from \(0.1544\pm0.0093\) to
\(0.1764\pm0.0255\). Its paired gain of \(0.0220\pm0.0344\) exceeds the
0.01 mean threshold, but only three of five differences are positive; the
95\% interval is [-0.0207, 0.0648]. This is a positive but unstable mean effect,
not evidence of uniformly harmful transfer. The four-of-five consistency
condition is not met, so RT-DETRv2 FAM is not promoted to Stage B.

YOLO26 FAM decreases best-checkpoint mean \mapfifty{} from
\(0.0507\pm0.0163\) to \(0.0273\pm0.0149\), with a paired difference of
\(-0.0234\pm0.0148\) and only one positive pair. All ten Base and FAM
runs complete fifty epochs. Final-epoch means fall to 0.001262 and 0.000092,
respectively, showing substantial degradation in both variants. The selected
best checkpoints are predominantly early because they are selected by the
declared validation rule over the full completed budget, not because the
comparison was restricted to the first epoch. YOLO26 FAM fails Stage A and
does not enter Stage B.

The negative YOLOv10 result is scientifically useful. It shows that the RT-DETR
gain depends on the surrounding representation, fusion interface, optimization
trajectory, or checkpoint regime and cannot be assumed from the local FAM
operation alone. RT-DETRv2 indicates that architectural similarity can coexist
with a favorable mean effect that remains insufficiently stable across paired
seeds. The low YOLO26 scores and final-epoch degradation limit attribution to
this dual-backbone port and training recipe, not the detector family in
general. None of these results shows that YOLOv10,
RT-DETRv2, YOLO26, or DINO is intrinsically unsuitable for RGB--thermal
detection.

\section{Overall interpretation}

The most defensible conclusion combines complementary experimental contrasts. The
historical locked study establishes that standard FAM improves Base
RT-DETR consistently on the internal benchmark and primarily raises recall.
The validation-led campaign does not confirm a replacement: RCRA has a
positive but unstable Stage-B mean effect, while P2, global upsampling,
zero-offset initialization, post-FAM gating, scalar residual calibration,
mixed consistency, and box-guided alignment fail their five-seed Stage-A
development rules.
Standard FAM therefore remains the selected model and the reference for
subsequent candidate comparisons.

That choice is strengthened by the previously unused acquisition campaign:
historical FAM beats historical Base in 10/10 acquisition--seed contrasts.
The detector-transfer comparisons do not establish a general FAM advantage.
RT-DETRv2 has a positive mean effect but wins only three of five seeds and
fails the same Stage-A consistency requirement. Within YOLO26, FAM loses to
Base in four of five paired seeds, while both
variants have very low final-epoch scores. Its different optimization budget
and evaluation pipeline preclude interpreting the absolute YOLO26 and RT-DETR
scores as a controlled cross-family ranking.

The paired modality audit strengthens one part of that conclusion and narrows
another. Fusion exceeds VIS on the same 708 frames in all five selected
checkpoints, so the thermal stream contributes useful information. Conversely,
masked IR in VIS coordinates is not interchangeable with native IR detection.
Reporting the coordinate contract is necessary to distinguish sensor
complementarity from missing-reference robustness.

At the same time, FAM is internally variable and computationally expensive.
One feature level can collapse without visible degradation in the headline
metric, a correction that makes the mechanism safer can make the detector
worse, and a more adaptive RCRA mechanism can improve the mean without being
stable across seeds. Architecture transfer and external scale mismatch further
show that an alignment-inspired inductive bias is not sufficient on its own.
Controlled synthetic perturbations add that higher absolute detection accuracy
and a more favorable signed relative-drop curve are distinct properties.

These results support a broader methodological point: evaluating multimodal
fusion should combine output metrics, coordinate-controlled modality tests,
paired replication, internal diagnostics, out-of-distribution stress, and
cost. Any one of these views alone would have produced an incomplete model
selection.

\section{Threats to validity}

\subsection{Internal validity}

Native CUDA execution is nondeterministic, particularly for deformable
operators. Pairing seeds controls intended stochastic sources but cannot pair
the microscopic CUDA perturbations of two architectures. Five independent
checkpoints estimate the resulting distribution but do not exhaust it.

The level ablations are interventions on already trained networks. Bypassed
representations are outside the training path and cannot be interpreted as the
performance of a newly trained P3-only or P3+P5 architecture. The P5 audit and
the choice of a bounded correction followed observation of the MtErie
checkpoint. Freezing one hypothesis before five new trainings reduces further
adaptation, but the correction remains post-hoc with respect to the internal
benchmark.

The standalone YOLOv10 evaluator differs slightly from the automatic
Ultralytics validator. The predefined tolerance is recorded as missed rather
than changed after inference. Comparisons among modalities use the same
standalone evaluator, while automatic VIS+IR results remain primary for the
training protocol.

The whole-sequence split, promotion thresholds, and Stage-B confirmation rule
were frozen before their respective comparisons. Nevertheless, the candidate
sequence was designed within the same study and after extensive prior
experimentation. Stage A is therefore an internal model-development procedure,
not an independent replication. The Modal Dropout coordinate-system ablation
uses only one predeclared screening seed and supports rejection of that recipe,
not a five-seed population estimate.

The zero-offset initialization, mixed-consistency, and Box-Guided comparisons
reuse the five completed standard-FAM Stage-A controls. Shared model,
training, data, and split settings were checked, and reconstructed
initialization hashes match for all five
seeds. For zero-offset initialization, the intervention additionally preserves
all unaffected parameters and random-number-generator state. Reconstructing
these initial states does not certify identical historical source snapshots
or environments. Mixed consistency adds student-path random draws that can
alter later Modal Dropout realizations. Seed pairing therefore evaluates the complete
training intervention rather than isolating the consistency loss under
identical intermediate tensors. Strict checkpoint loading and recorded replay
checks constrain execution and evaluation errors, but cannot remove these
attribution limits or native-CUDA divergence.

\subsection{Construct validity}

\mapfifty{} emphasizes detection at IoU 0.50 and is not a complete measure of
localization quality or SAR utility. Secondary COCO mAP, mAP@75, mAR, and error
metrics reduce dependence on one score but do not measure response time,
operator workload, alarm calibration, energy consumption, or consequences of
false alarms.

Recall--FPPI curves make the cost in false detections explicit but still count
boxes per frame, not distinct alarm events. Consecutive frames are correlated,
and tracking or temporal aggregation would change the operator-facing alarm
rate. These curves describe performance on evaluation sets that had already
been examined in earlier analyses. The shaded bands show variability
across five training seeds on these fixed sets, not uncertainty about
performance on new acquisitions. The confidence sweep does not establish
a confidence threshold for operational use, but it characterizes
the recall--false-positive trade-off.

Offset magnitude, feature cosine similarity, RMSE, smoothing, and PCA are
internal proxies. Without camera calibration or ground-truth correspondences,
they do not validate physical registration. PCA figures also depend on a
low-dimensional projection and are used illustratively.

Native-sensor and paired masked-modality AP operationalize different
constructs. The former changes the image canvas and ground truth with the
sensor, while the latter fixes a VIS-coordinate population and zeros channels.
The native-IR diagnostic is therefore evidence against branch collapse, not a
direct ranking against fusion. Likewise, \texttt{adapt\_ir2rgb} is a
resize-and-pad transformation and should not be interpreted as calibrated
registration.

The box-centre displacement used by Box-Guided P3 is a weak geometric proxy,
not dense pixel correspondence. Detection results alone do not establish
whether that target was learned consistently across seeds, let alone physical
registration or geometric generalization. Similarly, the signed relative drop
under synthetic translation and scaling confuses sensitivity with the
possibility of accidentally compensating a pre-existing mismatch.

\subsection{Statistical conclusion validity}

Five paired runs permit estimation of effect direction and variability but
produce wide intervals. Exact Wilcoxon tests have coarse resolution, and the
smallest attainable two-sided value for five nonzero signed pairs is 0.0625.
The reported tests are exploratory and are not corrected for the number of
contrasts. Checkpoints, rather than frames, boxes, diagnostic samples, or
feature cells, are the experimental units.

The Stage-A comparisons of zero-offset initialization, mixed consistency,
Box-Guided P3, RT-DETRv2, and YOLO26 likewise retain five paired checkpoints
per contrast. Their confidence intervals describe between-seed variability on
the fixed validation split, not
uncertainty across new acquisitions. Selection is governed by the predeclared
mean-gain and sign-consistency thresholds, not a post-hoc significance filter.

\subsection{External validity}

MtErie was repeatedly consulted during study development and is not a blind
holdout. Its locked and Stage-B comparisons are internally useful but may
overstate performance on a new acquisition. The official 273-frame FHL
validation remains too small and shifted for stable selection. The replacement
896-frame whole-video split is much more informative and removes same-video
adjacency, but FHL 0401/0402 and the remaining FHL training video still belong
to the same acquisition campaign.

Carnation 0025/0026 and FHL 0407/0408 were unused before the frozen
confirmation and contain substantially more paired frames, but they remain
internal WiSARD acquisitions. They strengthen acquisition-level replication
without demonstrating generalization outside WiSARD, to a different sensor
platform, or to the SAR population at large.

Carnation is one deliberately severe scale-mismatch case. It reveals a
specific boundary but cannot estimate performance across SAR domains. The
study uses one dataset, one main RT-DETR size, one YOLOv10 size, one
RT-DETRv2 configuration, one YOLO26 size, and one GPU for
latency. Different cameras, calibration, target scales, thermal conditions,
hardware, or training budgets may change the conclusions.

Deformable DETR and DINO results use an exploratory protocol. RT-DETRv2 and
YOLO26 follow the five-seed Stage-A decision procedure within the two-stage
framework, but neither satisfies the promotion rule and therefore neither
enters Stage B. Detector-specific recipes, budgets, and evaluators remain
different. Their Stage-A validation scores also have a different population
and checkpoint role from the full-data MtErie results. These tables cannot form
a quantitatively homogeneous architecture leaderboard with final RT-DETR and
YOLOv10.

\section{Future work}

\subsection{Independent confirmation of the model-selection protocol}

This thesis implements a train-derived, whole-video validation protocol and
shows that the five validation-selected FAM checkpoints outperform their
corresponding epoch-10 checkpoints on MtErie in all five Stage-A seeds. The
next priority is therefore not to create that split again, but to test
whether the rule transfers beyond the FHL acquisition campaign. A genuinely
independent paired sequence should be reserved before architecture development,
and the checkpoint metric, horizon, promotion thresholds, and allowed search
space should be fixed before observing it.

The RT-DETR ten-epoch design deliberately avoids early stopping. Future work
may compare fixed-budget checkpoint selection with early stopping, but only
under the same independent validation population and equal maximum compute.
This would help determine whether the late YOLOv10 degradation reflects
overfitting, schedule mismatch, or a detector-specific checkpoint problem.

\subsection{Improving the alignment architecture}

The negative and unstable results narrow the next useful hypotheses. Repeating
P2, global \(800\times800\) input, or another learning-rate search on the same
post-FAM gate would add little evidence. More targeted directions are:

\begin{enumerate}[label=\arabic*.]
  \item use camera calibration, dense paired cross-modal annotations, or a
        controlled acquisition rig with known physical geometry to supervise
        and evaluate alignment explicitly;
  \item train auxiliary modality-specific heads or coordinate-consistency
        losses so that native IR and paired VIS-coordinate robustness are not
        forced into one ambiguous objective;
  \item evaluate RCRA on an independent sequence and with more seeds before
        altering its descriptors or optimization, to distinguish a real
        conditional-alignment effect from training variance;
  \item revisit mixed native/paired consistency only with a new mechanism and
        an independent development population; the tested teacher--student
        formulation should not be retuned on the same FHL validation;
  \item reduce FAM or RCRA cost through channel compression or selective
        application to the most useful levels, followed by a complete inference
        benchmark rather than training-time measurements alone.
\end{enumerate}

These candidates should not be searched retrospectively on MtErie. A small
number of mechanistically justified variants should be frozen on an independent
development protocol and followed by paired replication and the same internal
audit used in this thesis.

\subsection{Broader evaluation and deployment}

Further work should include multiple external RGB--thermal acquisitions,
calibrated scale and translation experiments, additional missing-sensor
patterns, confidence calibration, and end-to-end latency including image
loading, preprocessing, transfer, postprocessing, and visualization. Deployment
experiments should report power and memory limits of the target UAV hardware
rather than extrapolating from a laptop GPU.

\section{Conclusions}

This thesis evaluated RGB--thermal fusion for UAV Search and Rescue from a
reliability-first perspective. Its study of the inherited RT-DETR Feature
Alignment Module combines paired multi-seed comparison, direct inspection of
learned transformations, temporal and whole-video validation audits,
acquisition-level stress testing, and configuration-matched candidate
confirmation.

The historical locked study shows that standard FAM improves Base RT-DETR in
all five paired seeds and consistently increases recall. The
validation-led campaign rejects P2, global upsampling, zero-offset
initialization, post-FAM gates, scalar residual control, mixed consistency,
and box-guided alignment under the common five-seed protocol. RCRA is the
only candidate to pass Stage A and learns the intended conditional mechanism, but
its full-data gain appears in only 3/5 seeds. Under the rule fixed before
Stage-B evaluation, standard FAM therefore remains the selected model.

The corrected paired modality audit shows that this FAM model benefits from
thermal information beyond VIS on the same 708 frames. It also exposes a sharp
coordinate-contract boundary: paired masked IR with VIS ground truth is very
weak, while native IR on native annotations remains functional. A one-seed
attempt to replace native dropout with the paired contract fails to resolve the
first condition and sharply reduces native-IR performance in that run.

The frozen confirmation on Carnation 0025/0026 and FHL 0407/0408 strengthens
the central architectural result: historical FAM exceeds Base in all ten
acquisition--seed comparisons, with a macro-average gain of 0.1257 \mapfifty{}.
At the same time, fusion does not uniformly exceed VIS on the matched Stage-B
checkpoints. The synthetic intervention also rejects the stronger claim that
FAM always degrades more gracefully under IR translation or rescaling.

The within-family transfer comparisons delimit the architectural result.
RT-DETRv2 FAM improves the best-checkpoint mean by 0.0220 but wins only three
of five seeds, so the Stage-A stability condition prevents promotion.
YOLO26 FAM has a lower best-checkpoint mean than its paired Base control
and wins only one of five seeds; both variants degrade strongly over the
fifty-epoch horizon. Full-budget replication and validation-selected checkpoints
are essential to interpret this result as a property of the tested recipe,
without extending it to all uses of FAM or all YOLO26 architectures.

The principal outcome is therefore both architectural and methodological. FAM
is the strongest confirmed RT-DETR configuration in this study, but it is
not a universal alignment solution and no new architecture is shown to improve
it reliably. Repeated runs, predeclared development and confirmation stages,
coordinate-controlled modality evaluation, internal diagnostics, frozen
acquisition confirmation, controlled stress cases, and computational cost are
necessary to determine whether an
apparent fusion gain is reliable, interpretable, and transferable.
